% Irodov, Arihant PDF pp. 17--19; original printed pp. 8--10.
% Problems 51--58: rolling, relative angular motion, cone and precession.

\exproblem{51}{8}
圆盘自初始圆心 $C=O$ 起沿 $x$ 正向平移并逆时针转动（原书图1.6）。
求瞬时转轴在平面内的轨迹 $y(x)$：
(a) 圆心以恒速 $v$ 平移，圆盘初始角速度为零并以恒定角加速度
$\beta$ 转动；(b) 圆心从静止起以恒定加速度 $w$ 平移，
圆盘以恒定角速度 $\omega$ 转动。
\exsolution
以圆心位置 $(x_C,0)$ 为参考，圆心瞬时速度
$\bm v_C=v_C\bm i$。平面内任一点 $P$
的速度为
\[
 \bm v_P=\bm v_C+\omega\bm k\times
 \bigl[(x_P-x_C)\bm i+y_P\bm j\bigr]
 =(v_C-\omega y_P)\bm i+\omega(x_P-x_C)\bm j.
\]
瞬时静止点满足 $\bm v_P=0$，故
\[
 x_P=x_C,\qquad y_P=v_C/\omega.
\]
(a) $x_C=vt$、$v_C=v$、$\omega=\beta t$，因此
\[
 \boxed{y(x)=\frac{v^2}{\beta x}}\quad(t>0),
\]
它是第一象限的一支双曲线。
(b) $x_C=wt^2/2$、$v_C=wt$，故
\[
 \boxed{y(x)=\frac{\sqrt{2wx}}{\omega}}\quad(t\ge0),
\]
其平方是关于 $x$ 的正比例关系。
\excheck 注意 (b) 中应带\emph{平方根}；
在 (a) 的 $t=0$ 时 $\omega=0$，圆盘作纯平移，
没有有限位置的瞬时转轴。
\begin{center}
\begin{tikzpicture}[scale=.75,>=Latex]
\draw[->] (-.2,0)--(5.2,0) node[right] {$x$};
\draw[->] (0,-.2)--(0,3.5) node[above] {$y$};
\draw[very thick,blue!65!black,domain=.7:5,samples=80] plot (\x,{2.1/\x});
\draw[very thick,red!70!black,domain=0:5,samples=80] plot (\x,{.92*sqrt(\x)});
\node[blue!65!black] at (1.15,3) {(a)};
\node[red!70!black] at (4.4,2.4) {(b)};
\end{tikzpicture}
\end{center}

\exproblem{52}{8--9}
半径为 $R=0.50\,\mathrm m$ 的轮子以恒定平移速度
$v=1.00\,\mathrm{m/s}$ 沿水平面无滑动滚动。
轮缘有一点 $A$。求：
(a) $A$ 的加速度大小与方向；
(b) 它相邻两次接触地面间走过的总路程。
\exsolution
无滑动滚动时角速度 $\omega=v/R$ 恒定，角加速度为零。
以轮心为参考，任一轮缘点只有匀速转动的向心加速度，
轮心自身加速度为零，所以
\[
 \boxed{|\bm a_A|=\omega^2R=\frac{v^2}{R}
 =2.00\,\mathrm{m/s^2}},
\]
且始终沿 $A$ 指向轮心。
将接触地面时刻取为 $\theta=0$，轮缘点的摆线参数方程为
\[
 x=R(\theta-\sin\theta),\quad
 y=R(1-\cos\theta),\qquad 0\le\theta\le2\pi.
\]
弧长元
\[
 \frac{\mathrm ds}{\mathrm d\theta}
 =R\sqrt{(1-\cos\theta)^2+\sin^2\theta}
 =2R\sin\frac\theta2.
\]
故整段轨迹长度
\[
 \boxed{s=\int_0^{2\pi}2R\sin(\theta/2)\,\mathrm d\theta
 =8R=4.00\,\mathrm m.}
\]
\excheck 这不是轮心的位移 $2\pi R$，而是轮缘点摆线的路程。

\exproblem{53}{9}
半径为 $R=10.0\,\mathrm{cm}$ 的小球沿斜面无滑动滚下，
球心沿斜面作匀加速运动，加速度 $w=2.50\,\mathrm{cm/s^2}$。
从静止出发 $t=2.00\,\mathrm s$ 后，按原书图1.7，
$A$ 为球面最上点、$B$ 为下坡方向的侧边点、
$O$ 为接触点。求三点的 (a) 速度；(b) 加速度。
\exsolution
取沿斜面向下为 $\bm i$、沿法线离开斜面为 $\bm j$。
在该时刻球心速率为
$u=wt=5.00\,\mathrm{cm/s}$；
角速度大小为 $\Omega=u/R$，角加速度大小为
$\dot\Omega=w/R$，转动方向为顺时针。
刚体点速度满足
$\bm v_P=u\bm i+\bm\Omega\times\bm r_{P/C}$，
因此
\[
 \boxed{\bm v_A=2u\bm i,\quad
 \bm v_B=u\bm i-u\bm j,\quad
 \bm v_O=\bm0.}
\]
速率分别是
$\boxed{v_A=10.0,\ v_B=5\sqrt2\simeq7.07,\ v_O=0\
 \mathrm{cm/s}}$。
加速度使用完整刚体公式
\[
 \bm a_P=w\bm i+
 \dot{\bm\Omega}\times\bm r_{P/C}
 +\bm\Omega\times(\bm\Omega\times\bm r_{P/C}).
\]
代入 $r_{A/C}=R\bm j,\ r_{B/C}=R\bm i,\
r_{O/C}=-R\bm j$ 可得
\[
 \boxed{\bm a_A=2w\bm i-\frac{u^2}{R}\bm j,\quad
 \bm a_B=\left(w-\frac{u^2}{R}\right)\bm i-w\bm j,\quad
 \bm a_O=\frac{u^2}{R}\bm j.}
\]
因 $u^2/R=2.50\,\mathrm{cm/s^2}$，
\[
 \boxed{a_A=\sqrt{5.00^2+2.50^2}\simeq5.59,\quad
 a_B=2.50,\quad a_O=2.50\quad\mathrm{cm/s^2}.}
\]
\excheck 接触点虽有瞬时速度零，却\emph{不}意味着其加速度也为零。
这里 $(\bm i,\bm j)$ 为随斜面方向选定的局部正交基。

\exproblem{54}{9}
半径 $r$ 的圆柱在水平面上以恒定速度无滑动滚动。
求圆柱表面顶点 $A$、最前侧点 $B$（位置参照原书图1.7）
所描摆线在该瞬间的曲率半径。
\exsolution
设圆心速率为 $u$、角速度 $\Omega=u/r$，
则两点速度和加速度的分量为
\[
 \bm v_A=(2u,0),\quad
 \bm a_A=(0,-u^2/r);
\qquad
 \bm v_B=(u,-u),\quad
 \bm a_B=(-u^2/r,0).
\]
平面轨迹的瞬时曲率半径可由
\[
 R_c=\frac{|\bm v|^3}{|\bm v\times\bm a|}
\]
计算，注意叉积取垂直平面的标量大小。因而
\[
 \boxed{R_A=\frac{(2u)^3}{2u^3/r}=4r,\qquad
 R_B=\frac{(\sqrt2u)^3}{u^3/r}=2\sqrt2\,r.}
\]
\excheck 两式均独立于滚动速率 $u$；
固定位置处的曲率半径仅由圆柱半径和点的位置决定。

\exproblem{55}{9}
两个刚体绕空间中固定、相交且互相垂直的轴转动，
各自角速度大小恒为 $\omega_1=3.0\,\mathrm{rad/s}$、
$\omega_2=4.0\,\mathrm{rad/s}$。
求刚体1相对刚体2的角速度与角加速度大小。
\exsolution
在任一瞬间，相对角速度作为空间矢量为
$\bm\omega_{1/2}=\bm\omega_1-\bm\omega_2$，
两轴垂直，故
\[
 \boxed{|\bm\omega_{1/2}|=
 \sqrt{\omega_1^2+\omega_2^2}=5.0\,\mathrm{rad/s}.}
\]
相对角加速度须在随刚体2旋转的坐标系中求导，
而不能简单对空间矢量差求导：
\[
 \left(\frac{\mathrm d\bm\omega_{1/2}}{\mathrm dt}\right)_2
 =\left(\frac{\mathrm d\bm\omega_{1/2}}{\mathrm dt}\right)_{\rm inertial}
  -\bm\omega_2\times\bm\omega_{1/2}.
\]
两空间固定轴转动速度矢量均不随惯性系时间改变，
于是
\[
 \bm\beta_{1/2}=-\bm\omega_2\times\bm\omega_1,\qquad
 \boxed{|\bm\beta_{1/2}|=\omega_1\omega_2
 =12\,\mathrm{rad/s^2}.}
\]
\excheck “相对角加速度”指相对旋转的角速度在旋转参照系中的导数；
若误在惯性系中求 $\mathrm d(\bm\omega_1-\bm\omega_2)/\mathrm dt$，
会错误得到零。

\exproblem{56}{9}
刚体角速度随时间变化为
$\bm\omega=at\,\bm i+bt^2\,\bm j$，给定
$a=0.50\,\mathrm{rad/s^2}$、
$b=0.060$（根据表达式量纲应为
$\mathrm{rad/s^3}$）。求在 $t=10.0\,\mathrm s$：
(a) 角速度、角加速度的大小；
(b) 两矢量的夹角。
\exsolution
直接求导：
\[
 \bm\beta(t)=\frac{\mathrm d\bm\omega}{\mathrm dt}
 =a\bm i+2bt\bm j.
\]
取 $t=10\,\mathrm s$，有
$\bm\omega=(5.0,6.0)\,\mathrm{rad/s}$、
$\bm\beta=(0.50,1.20)\,\mathrm{rad/s^2}$，故
\[
 \boxed{|\bm\omega|=\sqrt{61}\simeq7.81\,\mathrm{rad/s},\qquad
 |\bm\beta|=\sqrt{1.69}=1.30\,\mathrm{rad/s^2}.}
\]
夹角 $\gamma$ 满足
\[
 \cos\gamma=\frac{\bm\omega\cdot\bm\beta}
 {|\bm\omega||\bm\beta|}
 =\frac{5(0.5)+6(1.2)}{\sqrt{61}\cdot1.3}
 =\frac{9.7}{1.3\sqrt{61}},
\]
所以 $\boxed{\gamma\simeq17.2^\circ}$。
\excheck 原书本版将参数 $b$ 的单位也印为
$\mathrm{rad/s^2}$，但因它乘 $t^2$，应为
$\mathrm{rad/s^3}$；上述数值按公式进行一致量纲解释。

\exproblem{57}{9}
一圆锥在水平面上匀速无滑动滚动，圆锥半顶角
$\alpha=30^\circ$，底面半径 $R=5.0\,\mathrm{cm}$。
顶点 $O$ 通过铰链固定，锥轴水平，底面中心 $C$ 的速率
$v=10.0\,\mathrm{cm/s}$（图1.8）。
求：(a) 圆锥的瞬时角速度大小及其与竖直方向的夹角；
(b) 角加速度大小。
\exsolution
无滑动时接触水平面的母线是圆锥的瞬时转轴，位于锥轴下方，
与水平锥轴夹角为 $\alpha$，故与竖直方向的夹角为
$90^\circ-\alpha=60^\circ$。
锥顶到底面中心的距离是
$OC=R/\tan\alpha$，点 $C$ 到接触母线的垂直距离为
$d=OC\sin\alpha=R\cos\alpha$。
由 $v=\omega d$ 得
\[
 \boxed{\omega=\frac v{R\cos\alpha}
 \simeq2.31\,\mathrm{rad/s},\qquad
 \angle(\bm\omega,\text{竖直})=60^\circ.}
\]
圆锥轴以方位角速度
$\Omega=v/OC=v\tan\alpha/R$
绕竖直方向转动；因 $\omega$ 大小不变、方向跟随母线转动，
其水平分量大小为 $\omega\cos\alpha$，
故角加速度大小
\[
 \boxed{\beta=\Omega\omega\cos\alpha
 =\frac{v^2}{R^2}\tan\alpha
 \simeq2.31\,\mathrm{rad/s^2}.}
\]
\excheck $\omega$ 与 $\beta$ 的数量级相同是本题数值巧合，
但量纲不同，不能互相替代。

\exproblem{58}{10}
一个刚体相对水平轴 $AB$ 自转，其自转角速度
$\omega_0=0.50\,\mathrm{rad/s}$ 保持恒定。
从 $t=0$ 开始，轴 $AB$ 又以恒角加速度
$\beta_0=0.10\,\mathrm{rad/s^2}$
由静止绕竖直方向转动。
求 $t=3.5\,\mathrm s$ 时的角速度与角加速度大小。
\exsolution
记沿瞬时水平轴 $AB$ 的单位矢量为 $\bm e_h(t)$，
沿竖直方向的单位矢量为 $\bm k$。
水平轴的方位角速度大小为 $\Omega(t)=\beta_0t$，
于是刚体总瞬时角速度为
\[
 \bm\omega(t)=\omega_0\bm e_h(t)+\Omega(t)\bm k.
\]
由于两分量互相垂直，
\[
 \boxed{|\bm\omega|=\sqrt{\omega_0^2+\beta_0^2t^2}
 =\sqrt{0.25+0.1225}\simeq0.610\,\mathrm{rad/s}.}
\]
随轴旋转的单位矢量满足
$\dot{\bm e}_h=\Omega\,\bm k\times\bm e_h$，
故角加速度为
\[
 \bm\beta=\dot{\bm\omega}
 =\omega_0\Omega(\bm k\times\bm e_h)+\beta_0\bm k.
\]
两项也互相垂直，因此
\[
 \boxed{|\bm\beta|=\beta_0\sqrt{1+\omega_0^2t^2}
 =0.10\sqrt{1+0.25(3.5)^2}
 \simeq0.202\,\mathrm{rad/s^2}.}
\]
\excheck 不可只把 $\omega_0$ 与 $\Omega$ 视作两个静止方向的标量：
水平轴方向本身在改变，其方向导数
$\dot{\bm e}_h$ 正是额外的角加速度来源。
